% 2011 年天津卷文科第 3 题：程序框图（循环：|x|>3 时执行 x=|x-3| 再回判断）
% 来源 PDF：普通高考/2011/2011天津文.pdf，第 1 页（印面 707）第 1 栏。
% 重绘依据：原图为「开始 → 输入 x → 判断 |x|>3?」，是分支向下执行 x=|x-3| 并回到判断，
%   否分支向右再向下汇入 y=2^x，随后「输出 y」与「结束」。
\begin{tikzpicture}[scale=0.72, semithick, >=Stealth,
  block/.style={draw, minimum width=3.2cm, minimum height=0.6cm, align=center},
  dec/.style={draw, diamond, aspect=3.0, inner sep=1pt, align=center},
  arr/.style={->}]
  \node (start) [draw, rounded corners=6pt, minimum width=1.5cm, minimum height=0.5cm] {开始};
  \node (in)    [block, below=0.35cm of start] {输入 $x$};
  \node (dec)   [dec, below=0.5cm of in] {$|x|>3$?};
  \node (op)    [block, below=0.55cm of dec] {$x=|x-3|$};
  \node (pow)   [block, below=1.0cm of op] {$y=2^x$};
  \node (out)   [draw, trapezium, trapezium left angle=70, trapezium right angle=110,
                 minimum width=3.0cm, minimum height=0.6cm, below=0.45cm of pow] {输出 $y$};
  \node (end)   [draw, rounded corners=6pt, minimum width=1.5cm, minimum height=0.5cm,
                 below=0.4cm of out] {结束};
  \draw[arr] (start) -- (in);
  \draw[arr] (in) -- (dec);
  \draw[arr] (dec) -- node[right] {是} (op);
  \draw[arr] (op.west) -- ++(-1.6,0) |- (dec.west);
  \draw[arr] (dec.east) -- node[above] {否} ++(1.9,0) |- (pow.east);
  \draw[arr] (op) -- (pow);
  \draw[arr] (pow) -- (out);
  \draw[arr] (out) -- (end);
\end{tikzpicture}
